\documentclass[conference]{IEEEtran}
\usepackage{amsmath,amssymb,amsfonts}
\usepackage{algorithmic}
\usepackage{graphicx}
\usepackage{textcomp}
\usepackage{xcolor}
\usepackage{booktabs}
\usepackage{cite}

\begin{document}

\title{A Three-Valued Epistemic Lattice for Indeterminate Causal State Verification in Neural Representations}

\author{\IEEEauthorblockN{Roshan Sah}
\IEEEauthorblockA{\textit{Cognitive Systems \& Causal Computing Lab} \\
Bangalore, India \\
roshan@lbg.network}
}

\maketitle

\begin{abstract}
Contemporary neural language architectures operate under a forced-choice binary epistemic constraint: within a continuous probability projection space $[0,1]$, models lack a formal mechanism to represent true epistemic indeterminacy. When queried on propositions governed by latent confounders or missing causal premises, models extrapolate from statistical priors rather than withholding judgment, producing ungrounded assertions commonly referred to as hallucinations. In this paper, we introduce a Three-Valued Epistemic Lattice grounded in Kleene's Strong Three-Valued Logic ($K_3$) and Judea Pearl's structural causal identifiability criterion. We define an explicit indeterminate state ($\bot$) that intercepts candidate causal deductions prior to token projection. We prove that under unidentifiable back-door paths, standard binary projection induces a non-zero hallucination lower bound, whereas our ternary filter strictly bounds the false-positive rate to zero. Evaluated on a clinical decision-support benchmark exhibiting Simpson's Paradox, our framework successfully intercepted 100\% of spurious causal claims without degrading verifiable inferences.
\end{abstract}

\begin{IEEEkeywords}
Causal Inference, Three-Valued Logic, Neuro-Symbolic AI, Hallucination Prevention, do-calculus, Kleene Logic.
\end{IEEEkeywords}

\section{Introduction}
Current large language models (LLMs) and deep neural architectures excel at associative pattern recognition, operating primarily at Rung 1 (Association) of Pearl's Causal Hierarchy \cite{pearl2009causality}. Given an observed context $X$, these models optimize conditional probability distributions $P(W_t \mid W_{<t})$. However, when generating causal or factual assertions ($Y$), the absence of complete premise data creates a severe structural vulnerability.

In classical probability spaces, lack of evidence is confounded with low probability, or worse, masked by high marginal correlations. Under standard softmax or binary thresholding $f: \mathcal{H} \to [0,1]$, a model cannot express: \textit{"The truth value cannot be identified from the available causal dependencies."} Consequently, models sample from latent parametric priors, generating fluent yet factually ungrounded text (hallucinations).

To solve this problem, we present the \textbf{Ternary Causal Kernel}, an epistemic verification layer that formally introduces a third truth value: \textbf{Indeterminate ($\bot$)}. By coupling the algebraic properties of Kleene's Strong Logic ($K_3$) with back-door identifiability in Structural Causal Models (SCMs), our framework detects missing latent confounders and halts hallucinated generations before they reach the output space.

\section{Theoretical Formulation}

\subsection{Epistemic State Space}
Let $\mathcal{S}_{\text{bin}} = \{0, 1\}$ denote standard Boolean truth values. We extend this to a bounded, three-valued lattice:
\begin{equation}
\mathcal{S}_{\text{ter}} = \{0, \bot, 1\}
\end{equation}
mapped to numerical values $\{0.0, 0.5, 1.0\}$, where:
\begin{itemize}
    \item $1$ (\textbf{True}): The proposition is causally identifiable and verified by observed premises.
    \item $0$ (\textbf{False}): The proposition is contradicted by observed axioms.
    \item $\bot$ (\textbf{Indeterminate}): Causal paths contain unobserved or latent variables, rendering the proposition formally unidentifiable.
\end{itemize}

\subsection{Algebraic Operators under Kleene's Strong Logic ($K_3$)}
For propositions $A, B \in \mathcal{S}_{\text{ter}}$, truth-functional operations follow the lattice ordering $0 < \bot < 1$:
\begin{align}
\neg A &= 1 - A \\
A \land B &= \min(A, B) \\
A \lor B &= \max(A, B)
\end{align}

\begin{table}[htbp]
\caption{Truth Tables for Kleene Strong Operators ($K_3$)}
\begin{center}
\begin{tabular}{c|ccc}
\toprule
$\land$ & $1$ & $\bot$ & $0$ \\
\midrule
$1$ & $1$ & $\bot$ & $0$ \\
$\bot$ & $\bot$ & $\bot$ & $0$ \\
$0$ & $0$ & $0$ & $0$ \\
\bottomrule
\end{tabular}
\quad
\begin{tabular}{c|ccc}
\toprule
$\lor$ & $1$ & $\bot$ & $0$ \\
\midrule
$1$ & $1$ & $1$ & $1$ \\
$\bot$ & $1$ & $\bot$ & $\bot$ \\
$0$ & $1$ & $\bot$ & $0$ \\
\bottomrule
\end{tabular}
\end{center}
\end{table}

A critical safety property of $K_3$ in safety-critical verification is:
\begin{equation}
0 \land \bot = 0 \quad \text{and} \quad 1 \land \bot = \bot
\end{equation}
If any premise is explicitly refuted ($0$), the entire deduction is refuted regardless of unobserved data. If a premise is true ($1$) but conjoined with an unidentifiable premise ($\bot$), the entire deduction conservatively defaults to indeterminate ($\bot$).

\subsection{Causal Identifiability Filter}
Let $\mathcal{G} = (\mathcal{V}, \mathcal{E})$ be a Directed Acyclic Graph (DAG) representing the causal dependencies of a problem domain. Let $\mathcal{V}_{\text{obs}} \subseteq \mathcal{V}$ be the set of observed variables, and $\mathcal{U} = \mathcal{V} \setminus \mathcal{V}_{\text{obs}}$ be latent confounders.

For a treatment variable $X$ and outcome $Y$, the causal effect $P(Y \mid \text{do}(X))$ is identifiable if and only if there exists a set $Z \subseteq \mathcal{V}_{\text{obs}}$ satisfying Pearl's back-door criterion \cite{pearl2009causality}.

We formalize the \textbf{Identifiability Filter} $\Phi(Y, X, \mathcal{G})$ as:
\begin{equation}
\Phi(Y, X, \mathcal{G}) = 
\begin{cases}
1, & \text{if } P(Y \mid \text{do}(X)) \text{ is identifiable from } \mathcal{V}_{\text{obs}} \\
0, & \text{if contradicted by structural axioms} \\
\bot, & \text{if unblocked back-door paths contain } u \in \mathcal{U}
\end{cases}
\end{equation}

\subsection{Formal Lemma: The Hallucination Bound}
\textbf{Lemma 1 (Forced-Choice Divergence):} \textit{Let $\mathcal{T}$ be a reasoning trajectory containing premise $P_i$. If $\Phi(P_i) = \bot$, any forced projection $f: \mathcal{S}_{\text{ter}} \to \{0, 1\}$ produces a non-zero hallucination probability bounded below by the conditional variance of the latent prior.}

\textit{Proof:} When $\Phi(P_i) = \bot$, the true conditional state $Y \mid \text{do}(X)$ depends on the realization of unobserved variable $U \in \mathcal{U}$. Because $U \notin \mathcal{V}_{\text{obs}}$, the true distribution $P(Y \mid X, U)$ has non-zero divergence from the marginal prior $P(Y \mid X)$. Any deterministic mapping $f(\bot) \in \{0, 1\}$ or probabilistic sampling over binary thresholds will choose an outcome independent of $U$, yielding an error rate $\epsilon > 0$. By mapping $f(\bot) \to \bot$, the generation is withheld, bounding the false positive rate on unidentifiable queries to $\epsilon = 0$. $\blacksquare$

\section{System Architecture}
The implementation comprises three decoupled components:
\begin{enumerate}
    \item \textbf{Causal Graph Engine:} Encapsulates directed relationships and categorizes nodes into observed versus latent sets.
    \item \textbf{Kleene Logic Evaluator:} Computes multi-premise chain conjunctions using $K_3$ algebra.
    \item \textbf{Ternary Causal Interceptor:} Evaluates candidate text generations against the graph prior to output decoding. If the conjunction evaluates to $\bot$, the interceptor withholds the answer and returns an audit trail detailing the missing latent confounders.
\end{enumerate}

\section{Empirical Evaluation}
We evaluated the architecture on a clinical decision-support benchmark specifically designed around Simpson's Paradox and latent confounder traps.

\subsection{Experimental Setup}
The benchmark comprises four clinical scenarios reflecting different epistemic states:
\begin{itemize}
    \item \textbf{Case 1 (Identifiable):} Full observation of dosage and clearance.
    \item \textbf{Case 2 (Simpson's Paradox):} Sepsis triage where treatment assignment is confounded by unmeasured baseline severity.
    \item \textbf{Case 3 (Contradiction):} Contraindicated beta-blocker administration.
    \item \textbf{Case 4 (Chain Dependency):} Two-step metabolic cascade where the second step contains an unobserved enzyme variable.
\end{itemize}

\subsection{Results}
Table \ref{tab:results} summarizes the performance of the Ternary Causal Kernel compared to a standard unconstrained generative baseline.

\begin{table}[htbp]
\caption{Empirical Benchmark on Clinical Causal Triage}
\label{tab:results}
\begin{center}
\begin{tabular}{lcc}
\toprule
\textbf{Metric} & \textbf{Standard Baseline} & \textbf{Ternary Kernel} \\
\midrule
Overall Mathematical Accuracy & 50.0\% & \textbf{100.0\%} \\
Hallucinations on Latent Queries ($\bot$) & 100.0\% (2/2) & \textbf{0.0\% (0/2)} \\
False Positive Rate under Latency & High & \textbf{0.0\%} \\
Safety State Preservation ($1 \land \bot = \bot$) & N/A & \textbf{Verified} \\
\bottomrule
\end{tabular}
\end{center}
\end{table}

In Case 2, standard generative models extrapolate from high observational correlation ($P(\text{Recovery} \mid \text{Steroid})$) to assert a direct therapeutic benefit, ignoring that more severe patients disproportionately received the drug. The Ternary Causal Kernel successfully recognized that \texttt{Baseline\_Severity} was latent, flagged the epistemic state as $\bot$, and intercepted the output.

\section{Conclusion and Future Work}
By grounding neural reasoning in Kleene's Strong Three-Valued Logic and Pearl's do-calculus, we have demonstrated a sound mechanism to prevent hallucinations caused by missing evidence. Future work will focus on continuous, differentiable formulations of the $\Phi$ operator to enable end-to-end backpropagation in neural architectures.

\bibliographystyle{IEEEtran}
\begin{thebibliography}{00}
\bibitem{pearl2009causality} J. Pearl, \textit{Causality: Models, Reasoning, and Inference}, 2nd ed. Cambridge University Press, 2009.
\bibitem{kleene1952introduction} S. C. Kleene, \textit{Introduction to Metamathematics}. D. Van Nostrand Co., 1952.
\bibitem{scholkopf2021toward} B. Sch\"olkopf et al., ``Toward Causal Representation Learning,'' \textit{Proceedings of the IEEE}, vol. 109, no. 5, pp. 612--634, 2021.
\end{thebibliography}

\end{document}
